\section{Choosing a fixed-effect engine without changing the bank's question}
\label{exec20:fx-scope}
A bank comparing outreach across branches and months needs more than a fast
regression. It must know which account-periods remain in the calculation, whether
the treatment varies after the fixed effects are removed, what variation supports
uncertainty, and which customers an aggregated event-study effect describes.
These requirements motivate conditional reuse of the original \texttt{fixest}
engine through a typed R worker. The scope initially proposed here is linear
fixed-effect estimation; an IV formula, a nonlinear family or an event-study
expansion needs its own scientific and numerical qualification.

The inspected revision is commit \texttt{789f888416e9} of 26 June 2026, whose
\texttt{DESCRIPTION} declares version 0.14.2. This identifies the selected source,
without establishing the latest or a distributed release. Its R dependency is
at least 3.5.0, with Rcpp linking and imports including \texttt{numDeriv},
\texttt{nlme}, \texttt{sandwich}, \texttt{dreamerr} and \texttt{stringmagic}.
The declaration is GPL-3; the associated standard license text is available
separately. A root \texttt{LICENSE} file is absent from the inspected tree.
That absence does not erase the declaration or establish complete rights for
an assembled distribution. A separate process is an engineering choice and
does not settle distribution obligations. The selected change log records
covariance-class, IV-collinearity and degrees-of-freedom fixes; a workflow
specifies checks on three operating systems. These are maintenance evidence,
not observed passing checks or a support commitment.\citep{exec20_fixest_source}

\subsection{Projection, causal meaning and the retained population}
\label{exec20:fx-projection}
Let $Y$ be the vector of contribution margins before the incremental outreach
cost, $D$ the binary outreach indicator and $Z$ the full-rank matrix of branch
indicators. With equal observation weights, write
\begin{equation}
 M_Z=I-Z(Z^{\mathsf T}Z)^{-1}Z^{\mathsf T},\qquad
 \widehat\beta=(D^{\mathsf T}M_ZD)^{-1}D^{\mathsf T}M_ZY.
 \label{exec20:fx-within}
\end{equation}
The expression requires $D^{\mathsf T}M_ZD>0$. It is the coefficient obtained
by ordinary least squares on the explicit dummy matrix $[Z,D]$: minimizing
over branch intercepts first gives the residual objective
$\|M_Z(Y-D\beta)\|^2$, whose scalar normal equation yields
\eqref{exec20:fx-within}. For strictly positive weights $W$, the corresponding
argument applies to $W^{1/2}Y$, $W^{1/2}D$ and $W^{1/2}Z$. Consequently a
conversion that changes weights changes the projection being requested.

This algebra alone identifies no intervention effect. For example, a
constant-effect panel interpretation additionally needs a credible restriction
such as $\mathbb E[\varepsilon_{it}\mid D_{i1},\ldots,D_{iT},Z]=0$ in the
specified model, as well as consistent treatment, an appropriate interference
assumption and the required sampling conditions. Outreach triggered by an
unobserved deterioration in customer finances can violate that restriction.
An IV extension substitutes instrument relevance and substantive exclusion
and exogeneity assumptions; the formula syntax cannot supply them. Likewise,
adding time effects changes $Z$ and may remove the variation that previously
identified the projection. These are scientific choices that the proposed
business-to-model interface must state before submitting a regression.

Consider the following exact authored example, with two periods per branch.
The outcomes are hypothetical USD contributions per account-period.
\begin{center}
\begin{tabular}{c|cc|cc}
Branch & $D_1$ & $D_2$ & $Y_1$ & $Y_2$\\\hline
A & 0 & 1 & 10 & 12\\
B & 1 & 0 & 20 & 16\\
C & 0 & 1 & 30 & 36
\end{tabular}
\end{center}
Within-branch residualization gives $D^{\mathsf T}M_ZD=3/2$ and
$\widehat\beta=4$. Solving the explicit dummy normal equations gives the
same result, with intercepts $(9,16,31)$. Reversing input row order leaves
the coefficient and group scores unchanged when row identifiers travel with
their data. If instead treatment is one for every observation in B and zero
elsewhere, $D$ lies in the span of $Z$ and $M_ZD=0$. The coefficient is
unavailable; reporting zero would invent an answer to an unidentified
regression target. These are rational-arithmetic reference cases,
not results from running R.\citep{exec20_fixest_reference}

The inspected \texttt{feols} documentation describes recursive removal of
singletons and perfect-fit observations under its default setting. Its stated
slope invariance in the relevant removal setting does not imply invariant
inference or population reporting. The source also tracks collinear variables
and observation selections. For finite weights it rejects negative values
and removes zero weights; missing or infinite weights enter sample-selection
handling. Our initial adapter should prevalidate all weights and require an
explicit policy for exclusions, rather than rely on a message printed after
estimation. Every retained and removed row needs its original identifier and
reason. A factor level disappearing after exclusions can change rank, and a
dropped treatment must remain a typed unavailable result.\citep{exec20_fixest_source}

\subsection{A covariance request is part of the estimand service}
\label{exec20:fx-covariance}
The selected startup code sets the default covariance family to iid. Later
calls can read session options. Thus a formula containing branch fixed effects
does not by itself request branch-clustered errors. A reproducible worker
must set the covariance family, cluster identifiers, small-sample correction
and reference degrees of freedom explicitly. The source's \texttt{ssc}
interface separates $K$ adjustment, fixed-effect counting, cluster adjustment
and the degrees of freedom used for the Student reference distribution.
The default fixed-effect count is \texttt{nonnested}; \texttt{none},
\texttt{full} and optional exact counting express different conventions.
Its upstream tests describe an irregular two-way example with counts 6, 1,
8 and 7 under these choices. Those tests are useful future conformance
specifications and have not been executed here.\citep{exec20_fixest_source}

For the three-branch example let $\widetilde D=M_ZD$,
$\widehat u=M_ZY-\widetilde D\widehat\beta$ and
$s_g=\sum_{i\in g}\widetilde D_i\widehat u_i$. The uncorrected one-way
cluster sandwich for this scalar coefficient is
\begin{equation}
 \widehat V_0=\frac{\sum_gs_g^2}{(\widetilde D^{\mathsf T}\widetilde D)^2}
 =\frac{(-1)^2+0^2+1^2}{(3/2)^2}=\frac89.
 \label{exec20:fx-cluster}
\end{equation}
An explicitly chosen correction
$\widehat V_K=\widehat V_0\{(n-1)/(n-K)\}\{G/(G-1)\}$ gives
$4/3$ for $n=6,G=3,K=1$ and $10/3$ for $K=4$. The coefficient remains 4.
Choosing $G-1$ gives two reference degrees of freedom; it is distinct from
either variance scaling. This calculation isolates the choices an adapter
must preserve, without asserting that every native option or weighting law
uses these exact formulas. Three clusters provide no evidence of calibrated
coverage. More generally, dependence within a branch and between branches
must be matched to the uncertainty procedure.\citep{exec20_fixest_reference}

Numerical covariance repair adds another distinction. In the inspected code,
\texttt{vcov\_fix} defaults to true. A nonpositive eigenvalue triggers a
candidate repair; assignment of the repaired covariance occurs inside
\texttt{if(vcov\_fix)}. An adjacent comment says repair always occurs even
when the flag is false, but the executable branch has the narrower behavior.
The warning decision uses a maximum change greater than $10^{-8}$, so an
absence of a warning is not proof that nothing was changed. A proposed
adapter must preserve the requested policy and the original and repaired
matrices, or refuse to certify a repair state it cannot observe. Positive
definiteness is a numerical property; it does not validate a causal model,
cluster law or interval coverage.\citep{exec20_fixest_source}

\subsection{Event-time support and the meaning of an aggregate}
\label{exec20:fx-eventtime}
Suppose branches adopt a servicing policy in different calendar months.
The cohort is the adoption month and relative time is the observed month
minus that cohort. The inspected \texttt{sunab} documentation warns that,
when relative time is inferred from calendar values, a cohort absent from
the supplied period support is treated as never treated. Its example uses
cohort 1994 and observed years 1995--2005: an already-treated cohort can be
misclassified as never treated. The proposed conversion must therefore
construct treatment status from the business history before invoking the
helper. It must distinguish never treated, treated before observation and
not yet treated, and preserve reference periods, binning and unsupported
cells. A valid design further needs the substantive comparison assumptions,
including appropriate parallel trends and no anticipation; a matrix expansion
does not verify them.\citep{exec20_fixest_source}

The selected aggregation routine constructs shares from the signs of design
columns, using observation weights when requested, normalizes those shares
and averages the corresponding coefficients. It calculates uncertainty from
the full coefficient covariance. For a fixed vector of shares $a$ and
coefficient vector $\widehat b$, the target and conditional variance are
\begin{equation}
 \widehat\theta=a^{\mathsf T}\widehat b,\qquad
 \widehat{\operatorname{Var}}(\widehat\theta)=a^{\mathsf T}\widehat V a.
 \label{exec20:fx-aggregate}
\end{equation}
The second identity follows by expanding the covariance of a linear
combination. If shares themselves are estimated for a different population
target, their uncertainty needs an additional argument; treating the observed
shares as fixed does not settle that question.\citep{exec20_fixest_source}

For an authored contrast take $\widehat b=(2,8)^{\mathsf T}$,
$a=(3/4,1/4)^{\mathsf T}$ and
$\widehat V=\left(\begin{smallmatrix}4&1\\1&9\end{smallmatrix}\right)$.
Then $\widehat\theta=7/2$ and the variance is $51/16$. Dropping the covariance
cross term gives $45/16$; averaging the two standard errors gives $9/4$,
which is not $\sqrt{51}/4$. Equal coefficient shares would instead give an
effect of 5. At a stipulated incremental cost of USD 4 per account, these two
different targets imply net margins of USD $-0.50$ and USD 1. Neither is an
observed bank return; the contrast shows why silently changing aggregation
weights can reverse a proposal's recommendation.\citep{exec20_fixest_reference}

\subsection{A conditional reuse choice and a public qualification path}
\label{exec20:fx-choice}
The original R engine is the preferred proposal route because its sample,
design and covariance behavior can be kept together behind one typed boundary.
\texttt{feols}, \texttt{feglm} and \texttt{femlm} are distinct model surfaces
within that package; qualifying one linear profile does not qualify all of
them. \texttt{pyfixest} remains a separately assessed implementation.
An explicit-dummy reference is attractive for small transparent examples,
while a direct port of numerical routines would also have to reproduce
formula coding, row selection, fixed-effect rank, covariance defaults,
aggregation and failures. No measured cost or target-platform requirement
currently justifies that duplication.\citep{exec20_fixest_source,exec20_fixest_reference}

The proposed request carries the immutable data snapshot, units, row keys,
calendar, design, target population and uncertainty choices. The response
carries coefficient identities, covariance, retained sample, exclusions,
rank and convergence diagnostics, aggregation shares and any repair status.
Library, CLI and MCP surfaces must preserve those meanings. A missing target,
invalid request, unsupported model, numerical failure and cancelled job
remain different outcomes. An agent may propose a reviewed design, but an
initial tool should accept bounded formula structures rather than arbitrary
generated R expressions. These are proposed contracts; no such fixest worker
has been admitted by this source review.

Public and generated panels can test the merit of that design before a bank
experiment. Start with exact dummy/within agreement and an absorbed treatment,
then test row permutations, zero/missing weights, singleton cascades,
near-collinearity, covariance defaults and correction choices. Event-time
cases must include adoption outside the observed calendar and aggregation
with nonzero covariance. Native runs should compare outputs, masks and typed
failures against the authored references and independently inspected upstream
specifications. A complete pinned distribution, compiler/runtime environment
and dependency rights review precede executable admission. Subsequent public
data studies may assess numerical stability and workflow value; they cannot
establish a private bank's assignment mechanism or actual return.

The scientific owner must own identifying assumptions and dependence; the
data owner must own calendars, weights and exclusions; the adapter maintainer
must own conversions and upgrade conformance. Operational costs to measure
include worker startup, serialization, memory, rank and covariance computation,
dependency maintenance and diagnostics. This conditional proposal selection
can be completed now while native qualification and deployment remain open.
It converts a broad intention to reuse panel software into a testable design
without treating an available coefficient as a validated causal conclusion.
